%%==========================================================================
%%  saus-exam-starter.tex
%%  A question paper for THE SHAIKH AYAZ UNIVERSITY, SHIKARPUR.
%%  Copy this file, fill in the details and write the questions.
%%  Compile with:  latexmk saus-exam-starter.tex   (LuaLaTeX via .latexmkrc)
%%
%%  Two runs settle the marks table and "Page 1 of n"; latexmk does that.
%%==========================================================================
% type       = mid (default) | final | sessional | quiz | makeup
%              | supplementary | practical      -- names the paper
% solutions  = 0 (default) | 1                  -- 1 prints the answer key
% clo        = 1 (default) | 0                  -- CLO and level in the margin
% student    = 1 (default) | 0                  -- roll number strip
% rollhead   = 1 (default) | 0                  -- "Roll No." on later pages
% fontsize   = 12pt (default) | 11pt | 10pt
\documentclass[type=mid]{saus-exam}

%% ---- The paper -------------------------------------------------------------
\sausdepartment{Department of Computer Science}
\examcourse{CSC-201}{Data Structures and Algorithms}
\examprogramme{BS Computer Science}
\examsemester{III}
\examsession{Spring 2026}
\examdate{14 April 2026}
\examduration{90 minutes}
\exammarks{30}
% \examcredits{3 (3-0)}                 % optional
% \examteacher{Name of the teacher}     % optional
% \examinfo{Venue}{Examination Hall 2}  % any further detail, in pairs
% \examtitle{Pre-Board Examination}     % overrides the wording of type=

\begin{document}
\makeexamhead

\begin{instructions}
  \item Attempt all questions in the answer book provided.
  \item Marks for each question are shown in the margin.
  \item Programmable calculators and mobile telephones are not allowed.
\end{instructions}

%% ---- Questions -------------------------------------------------------------
%% \question takes marks, clo, bloom and topic; clo, bloom and topic feed the
%% distribution table at the end. lines=n or space=40mm leaves room to write.
\examsection[10 marks]{Section A: Short questions}

\question[marks=5, clo=1, bloom=Understand, topic=Complexity]
Write the question here.

\question[marks=5, clo=2, bloom=Apply, topic=Stacks, lines=6]
A question answered on the paper itself, with six ruled lines after it.

\examsection[20 marks]{Section B: Long questions}

\question[marks=20, clo=3, bloom=Analyse, topic=Sorting]
A question in parts. The marks on the question cover its parts, so the
parts' own marks are not counted twice.
\qpart[marks=12]
The first part.
\qpart[marks=8]
The second part.

% \question[marks=3, bonus, clo=3] Bonus questions are left out of the total.

\examend

%% ---- For the department's file ---------------------------------------------
%% Delete this page before printing the candidates' copies.
\clearpage
\sausmarkstable
\end{document}
